\chapter{Logarithme, puissances et modèles multiplicatifs}
\section*{Exercice 16.1 -- Domaine et produit}
Le domaine est $x>2$. L'équation devient $x(x-2)=3$, soit $(x-3)(x+1)=0$. Seul $x=3$ est admissible ; $-1$ vient d'une équation algébrique sur un domaine trop large.

\section*{Exercice 16.2 -- Seuil}
En prenant les logarithmes et divisant par $\ln0{,}8<0$, $n>\ln0{,}1/\ln0{,}8\simeq10{,}32$. Le premier entier est $11$. On peut contrôler $0{,}8^{10}>0{,}1$ et $0{,}8^{11}<0{,}1$.

\section*{Exercice 16.3 -- Limite en zéro}
Poser $t=-\ln x\to+\infty$ : l'expression vaut $-te^{-2t}\to0$ par croissance comparée. Le signe reste négatif près de zéro, ce qui est compatible avec une limite nulle.

\section*{Exercice 16.4 -- Puissance réelle}
$x^{\sqrt2}=e^{\sqrt2\ln x}$. La composition donne une dérivée $e^{\sqrt2\ln x}\sqrt2/x=\sqrt2 x^{\sqrt2-1}$. La base positive est indispensable pour cette définition.

\section*{Exercice 16.5 -- Logarithme d’une somme}
Le domaine est $x>0$. Le membre de droite vaut $\ln(2x)$, donc par injectivité $x+1=2x$, d'où $x=1$. On ne sépare jamais $\ln(x+1)$ en deux logarithmes.

\section*{Exercice 16.6 -- Une limite exponentielle}
Le logarithme de l'expression vaut $n\ln(1+2/n)=2\,[\ln(1+2/n)/(2/n)]\to2$ par la dérivée de $\ln$ en $1$. La continuité de l'exponentielle donne la limite $e^2$.
